\begin{abstract}
We show how input symmetry can neutralize fresh shuffling in \scheme{}, an NDSS 2026 defense that combines secret linear transformations, additive markers, and independently shuffled KV-cache blocks. A repeated-token text prompt makes the first-layer Value rows identical. Shuffling then changes only the marker's position: protected blocks occupy an orbit of at most $b$ states, despite $b!$ possible permutations. Rank-one differences between these states reveal an equivalent row demixer and the implementation's structured column transform. The recovered parameters transfer to other requests protected by the same secrets, exposing token multisets without recovering within-block order. We characterize this orbit, establish identification under explicit nondegeneracy conditions, and derive a query bound. Against the official implementation, the attack recovers 263,029 of 263,904 head--token evaluation instances across ten models and twelve configurations. All fifteen additional model--key epochs calibrate successfully, with model-level mean recovery between 99.34\% and 100\%. A limited study of 30 synthetic closed-set values under one key yields 77.36\% expected top-1 identification under uniform tie-breaking; this is not a population-level privacy estimate. Request-specific refresh of the row transform and marker blocks the tested cross-request transfer, without establishing a complete secure repair. These results show why fresh permutations can leave reused secrets exposed when input symmetry reduces their observable effect to a small marker orbit.
\end{abstract}
